\chapter{Experiments}

\section{Introduction}

This chapter details how the experiments were conducted. Chapter~4 described the methodology and evaluation framework; however, it was the focus of this chapter to detail the actual experiment process: software setup, model/service startup, notebook order of execution, question scenarios, parameter settings, logging of results, and output validation.

These experiments have been divided into two phases (Local and HPC) that both run the same six \ac{RAG} pipeline configurations and the same five scenarios. Local is used to determine baseline performance against which all six \ac{RAG} pipeline variations will be compared. \ac{HPC} runs these same six pipeline configurations and the same five scenarios within the Cyclone cluster environment. Although local and HPC workflows do not represent a pure hardware comparison due to differences in the serving stack and backend models between environments.

\section{Experimental Setup}

\subsection{Software Environment}

The experiments were run in a set of Python Jupyter notebooks. Each experiment was executed within its own virtual environment containing all required packages. All the notebooks were run from their respective project notebook directory and contained an identical flow through the various sections described below: Path configuration; Service configuration; Corpus load; Split documents into training sets and test sets; Indexing; Build pipelines; Run scenarios; Log results; Validate results

The executed implementation notebooks are:

\begin{itemize}
    \item \texttt{rag\_\allowbreak basic.ipynb};
    \item \texttt{rag\_\allowbreak vectordb.ipynb};
    \item \texttt{rag\_\allowbreak vectordb\_\allowbreak tuned.ipynb};
    \item \texttt{rag\_\allowbreak reranker.ipynb};
    \item \texttt{rag\_\allowbreak hybrid.ipynb}; and
    \item \texttt{rag\_\allowbreak hybrid\_\&\_\allowbreak rerank.ipynb}.
\end{itemize}

All notebooks followed a similar general process of execution. That process included importing libraries, configuring paths and services, loading the document corpus, splitting/indexing documents, building the RAG pipeline, executing scenario questions, and saving results. This standardized structure reduces unwanted variation between implementations and makes the pipeline-level comparison more controlled.

\subsection{Models and Services}

The local process utilizes Ollama for model serving. The local generative model utilized is \texttt{llama3.1:8b} and the local embeddings model utilized is \texttt{mxbai-embed-large}. Both of these models are open-source and may be downloaded via Ollama or model repository sources (e.g., Hugging Face) subject to download restrictions on each model. Each notebook that runs reranked vectors locally utilizing \ac{MiniLM} cross encoder (sentence-transformers).

Each of the vector-database, reranking, hybrid, tuned-vector-database, and hybrid-reranking notebooks run prior to launch of Milvus, utilize the standalone script \textit{embed.sh} (project script) to initiate the Milvus service. The in-memory baseline is simple to implement, requires no Milvus services to begin with; hence it serves as the simplest baseline implementation.

In contrast, the \acf{HPC} workflow utilizes \ac{vLLM}-based model serving. Hence, the generation service utilized was \texttt{meta-llama/Llama-3.1-8B-Instruct}, the embedding service utilized was BAAI/bge-m3, and the reranking service utilized was \texttt{BAAI/bge-reranker-base}. Therefore, both workflows are based upon the same pipeline structure and timing category, however they are different model-serving deployments. The \ac{HPC} comparison should be viewed as a comparison at the workflow level and not strictly as an isolated hardware performance metric.

\subsection{Input Corpus}

All experiments make use of the same OrchestrAI synthetic enterprise corpus. The corpus includes 54 DOCX documents divided by departments. Each notebook loads each document recursively and converts it into a Haystack document format; it also keeps track of source information about each document (file name and department).

Since all notebooks have access to the same set of documents, the primary cause for differences in answers is the way in which they retrieve and rank those documents. The same corpus is duplicated on the \ac{HPC} system used for the second phase and retains the directory structure and file names.

\subsection{Configuration and Parameters}

The key configurable options include the model types, query types (e.g. "what," "where"), total number of retrieved documents, chunking configurations and logging schema. In most cases a total of ten retrieved document-chunks are passed to the language model for evaluation. For the reranked versions, the retriever first returns twenty candidates, and the cross-encoder reranker filters them down to the top ten.

All of the chunking configurations differ from each other only when they were used as part of an experimental design. The baseline used 120 word chunks that had 20 words of overlap. All of the Milvus-supported variants, the hybrid version and the hybrid-reranking variant all used 80 word chunks with 20 words of overlap. The only chunk size used by the vector database tuning notebook was 200 words with 40 words of overlap. The rest of the vector-database configuration remained unchanged.

\begin{itemize}
    \item \textbf{Top 10 retrieval:} basic, vector database, tuned vector database, and hybrid.
    \item \textbf{Top 20 retrieval followed by top 10 reranking:} reranking and hybrid with reranking.
    \item \textbf{Chunking test:} tuned vector database changes only chunk length and overlap.
\end{itemize}

\section{Deployment Stages}

\subsection{Local Execution Stage}
The local execution stage provides the baseline experimental stage of the project. The primary goal of this stage is to validate that each notebook successfully runs on the author’s computer and to provide comparative timing, retrieval and answer quality outputs.

The notebooks were run in the same order as follows:

\begin{enumerate}
    \item baseline in-memory \ac{RAG};
    \item Milvus vector-database \ac{RAG};
    \item tuned Milvus vector-database \ac{RAG};
    \item dense retrieval with reranking;
    \item hybrid retrieval; and
    \item hybrid retrieval with reranking.
\end{enumerate}

This order represents increasing pipeline complexity, starting with the in-memory baseline and then adding persistent vector storage, chunking variation, reranking, hybrid retrieval, and hybrid retrieval with reranking. All outputs generated during each run are recorded into a single \texttt{results.csv} file for comparison purposes.

\subsection{HPC Execution Stage}

The \ac{HPC} execution stage involves the same pipeline sequence and scenario questions as for the local execution stage, with the difference being that model serving is done via \ac{vLLM}, with generation, embedding, reranking services, and that execution is through Slurm-managed jobs.

The \ac{HPC} experiments were run on the Cyclone \ac{HPC} facility, which has GPU nodes with: \(2\times\) 20-core Intel Xeon Gold 6248 CPUs, \(4\times\) NVIDIA V100 accelerators, and 192GB of memory per node. The service launch scripts submit separate Slurm jobs for the generation model, the embedding model, and the reranking model. Each service job requests one node, one GPU, one task per node, and ten CPU cores.

The \ac{vLLM} generation service uses \texttt{meta-llama/Llama-3.1-8B-Instruct}; the embedding service uses \texttt{BAAI/bge-m3}; the reranking service uses \texttt{BAAI/bge-reranker-base}. Each service then writes its host, port, URL, API key, and model name into an environment file under \texttt{/nvme/scratch/\$\{USER\}}. Finally, these environment files are used by the \ac{HPC} notebooks to connect to the right model-serving endpoints.

The purpose of this stage is to compare application-level timing behaviour across environments. Queue time, service startup time, and data-transfer time are deployment issues worth considering, but we won’t treat them as separate measured variables in this thesis. We’ll be focusing on timing comparisons of indexing time, retrieval time, generation time, and total runtime.

\section{Experiment Scenarios}

\subsection{Scenario-Based Query Execution}
Five predetermined question-answering scenarios serve as the framework for the experiments. Different retrieval and reasoning behaviours are tested in each case. A straight factual lookup is the first scenario. The latter situations call for cross-document or procedural reasoning.

The selection of five scenarios provides a focused but diverse comparison set rather than a statistically exhaustive benchmark. The scenarios were chosen to represent common enterprise question types: single-fact lookup, procedural checklist retrieval, cross-team operational reasoning, policy exception handling, and multi-hop customer-risk analysis. Using the same five scenarios across all six implementations keeps the evaluation controlled while still testing whether each pipeline handles different retrieval and reasoning demands.

The five scenario questions are:

\begin{enumerate}
    \item \textbf{\textit{What is the first-response SLA for a Sev-1 support incident?}} This is the clean baseline scenario. It checks whether the system can retrieve and answer a single operational fact accurately.
        \begin{itemize}
            \item Retrieval type: single-document
            \item Main test: exact fact retrieval
        \end{itemize}
    \item \textbf{\textit{What steps must be completed before a new employee is considered ready for day one at OrchestrAI?}} This scenario is the broadest scenario because it requires the system to connect support, product, logistics, and commercial-risk information into one operational response.
        \begin{itemize}
            \item Retrieval type: checklist or procedure retrieval
            \item Main test: whether the system can recover and summarize a structured process
        \end{itemize}
     \item \textbf{\textit{A new engineer starts on Monday, but their laptop has not arrived and their required access is still pending. Which teams are involved and what should happen next?}} This scenario requires combining information from more than one team to answer a realistic workplace problem.
        \begin{itemize}
            \item Retrieval type: multi-document, cross-team retrieval
            \item Main test: whether the system can synthesize HR, IT, and access-related responsibilities
        \end{itemize}
     \item \textbf{\textit{Can OrchestrAI make an urgent hardware purchase without following the normal PO workflow if a warehouse outage is at risk?}} This is the most demanding scenario in the notebook because it requires cross-functional reasoning across several business areas.
        \begin{itemize}
            \item Retrieval type: policy retrieval with exception logic
            \item Main test: whether the system can identify when a standard workflow may be bypassed
        \end{itemize}
     \item \textbf{\textit{A customer renewal is at risk because support cases are recurring, a requested feature is still unavailable, and replacement hardware is delayed. What should OrchestrAI do next, and which teams must be involved?}} This scenario is slightly more complex because the answer is not a single fact but a short operational process.
        \begin{itemize}
            \item Retrieval type: multi-hop, multi-document retrieval
            \item Main test: whether the system can combine support, product, logistics, and commercial risk signals into one grounded answer
        \end{itemize}
\end{enumerate}

The same questions are used for all implementations. Each implementation uses its own scenario identifier prefix so that rows remain easy to distinguish in the results file. For example, the baseline uses \texttt{b1}--\texttt{b5}, the vector-database run uses \texttt{v1}--\texttt{v5}, and the hybrid-reranking run uses \texttt{hr1}--\texttt{hr5}.

\subsection{Baseline RAG Experiment}

In this baseline experiment, the five scenarios are executed with the in-memory \ac{RAG} pipeline. The DOCX chunks are loaded in memory and indexed into an in-memory document store. For each of the five questions, the system retrieves the top 10 dense match candidates. This experiment serves as a reference point for all other experiments presented in this thesis.

\subsection{Vector Database RAG Experiment}

This experiment executes the same five scenarios, but now on a persistent vector database (Milvus). For each of the five questions, the system searches for the top 10 dense match candidates. This experiment examines the impact of persistently storing vectors in a database and whether the quality of the retrieved matches stays similar to those produced by the baseline experiment.

\subsection{Tuned - Vector Database RAG Experiment}

The following experiments execute the same five scenarios by running the Milvus vector-database pipeline at a larger chunk size with higher overlaps. Chunk size will impact how much context is included for each retrievable unit. To remove any other influences from the chunk-size on the retrievals (i.e., to simply determine the effect of the chunk-size), all else constant are kept; that is, the retriever, model, prompt and top-k values are identical to those used within the vector-database configuration. Thus, these experiments are referred as the "tuned" variant of the vector-database experiments.

\subsection{Reranking Experiment}

For this reranking experiment, the five scenarios are executed using dense retrieval (Milvus) followed by cross-encoder reranking. First, the retriever delivers 20 candidates and then the reranker selects the best 10 candidates for prompting. This experiment assesses the trade-off between additional latency during the retrieval stage versus improved ranking order.

\subsection{Hybrid Retrieval Experiment}

The hybrid retrieval experiment combines dense semantic retrieval with a sparse BM25-style signal generated inside Milvus. It retrieves 10 documents for each question. This experiment evaluates whether combining semantic and lexical signals improves robustness for questions that contain exact operational terms such as severity labels, team names, procurement terms, and workflow names.

\subsection{Hybrid Retrieval with Reranking Experiment}

In the hybrid reranking experiment, cross-encoder reranking is applied after hybrid candidate retrieval. The hybrid retriever provides 20 candidates and the reranker chooses the best 10. This is the most complex pipeline currently implemented. This experiment tests whether hybrid recall combined with reranked precision generates better answers than previous implementations at a higher computational cost.

\section{Benchmarking and Data Collection}

All outputs will be logged into the project’s results logger. Before each run of the notebooks, rows for the respective implementations will be deleted so that multiple runs will replace previous data instead of producing duplicate data points. Every row will contain the name of the implementation, ID of the scenario, question, answer provided by the system, source documents, departmental information based on available documentation, available context, Number of documents available, top-k parameters, time it took to index the documents (one-time), time it takes to retrieve answers (query-time) and time it takes to generate an answer (query-time).

In the last local vs HPC comparison of timing, these columns will be utilized to create a total time run and speed up analysis for each pipeline. Once all the scenarios have been ran with their respective pipelines, the results file will be used to create the manual evaluation worksheet. Using evidence from the DOCX corpus, scores and supporting documents will be populated into the manual evaluation worksheet. The deterministic evaluation script will calculate the answer quality metrics and retrieval quality metrics and create per-scenario as well as summary CSV files. The visualization notebook will utilize these processed outputs to generate plots of answer quality, retrieval quality, scenario-based analysis, timing-based analysis and local vs. HPC comparison for chapter~6

\section{Data Validation}

The recorded output is validated as complete prior to analysis. The validation entails verification that there are five rows in each of the implementations, that the generated answers were not null, that there were sources that were retrieved, and that the time fields were available. Additionally, all retrieved contexts are manually examined when developing the truth table (ground-truth sheet), since the quality score of the answer cannot be assumed; it must have some form of evidence based upon a subset of the data within the corpus.

For the \ac{HPC} experiments, validation focuses on whether the same six pipeline variants and five scenarios are present in the timing outputs and whether the timing fields required for local versus \ac{HPC} comparison are available.

\section{Summary}

This chapter outlined how the experimental setup and implementation were carried out for all of the six developed \ac{RAG} pipelines. It also discussed the environment where the software was executed, the model-serving setup used to deploy models, the input corpus, the parameters that were set during execution, the sequence in which notebooks are run, the phases of deploying an application, the types of question asked by each scenario, how data was collected from each experiment, and what steps were taken to validate the results of experiments conducted with this pipeline.

In particular, the local phase is compared to provide a point-of-reference for comparing retrieval methods and the \ac{HPC} phase uses the same structure and scenarios in a Cyclone cluster environment (using \ac{vLLM}-based services) to apply the same pipeline structures. Therefore, since both environments utilize different model-serving back-ends and some different back-end models, when comparisons are made across environments they should be viewed as a comparison of workflow execution times versus strictly as a hardware comparison. Data generated in this way is then analyzed in chapter ~6